%----------------------------------------------------------------------------------------
%   Доорх хэсгийг өөрчлөх шаардлагагүй
%----------------------------------------------------------------------------------------
%!TEX TS-program = xelatex
%!TEX encoding = UTF-8 Unicode
\documentclass[12pt,A4]{report}

\usepackage{fontspec,xltxtra,xunicode}
\setmainfont[Ligatures=TeX]{Times New Roman}
\setsansfont{Arial}

% \usepackage[utf8x]{inputenc}
% \usepackage[mongolian]{babel}
%\usepackage{natbib}
\usepackage{geometry}
%\usepackage{fancyheadings} fancyheadings is obsolete: replaced by fancyhdr. JL
\usepackage{fancyhdr}
\usepackage{float}
\usepackage{afterpage}
\usepackage{graphicx}
\usepackage{amsmath,amssymb,amsbsy}
\usepackage{dcolumn,array}
\usepackage{tocloft}
\usepackage{dics}
\usepackage{nomencl}
\usepackage{upgreek}
\newcommand{\argmin}{\arg\!\min}
\usepackage{mathtools}
\usepackage[hidelinks]{hyperref}

\usepackage{algorithm}
\usepackage{algpseudocode}

\usepackage{listings}
\DeclarePairedDelimiter\abs{\lvert}{\rvert}%
\makeatletter
\usepackage{caption}
\captionsetup[table]{belowskip=0.5pt}
\usepackage{subfiles}

\usepackage{listings}
\renewcommand{\lstlistingname}{Код}
\renewcommand{\lstlistlistingname}{\lstlistingname ын жагсаалт}

\usepackage{color}
\definecolor{codegreen}{rgb}{0,0.6,0}
\definecolor{codegray}{rgb}{0.5,0.5,0.5}
\definecolor{codepurple}{rgb}{0.58,0,0.82}
\definecolor{backcolour}{rgb}{0.99,0.99,0.99}
 
\lstdefinestyle{mystyle}{
    basicstyle=\ttfamily\small,
    backgroundcolor=\color{backcolour},   
    commentstyle=\color{codegreen},
    keywordstyle=\color{magenta},
    numberstyle=\tiny\color{codegray},
    stringstyle=\color{codepurple},
    %basicstyle=\footnotesize,
    breakatwhitespace=false,         
    breaklines=true,                 
    captionpos=b,                    
    keepspaces=false,                 
    numbers=left,                    
    numbersep=10pt,                  
    showspaces=false,                
    showstringspaces=true,
    showtabs=false,                  
    tabsize=2
}
 
\lstset{style=mystyle, label=DescriptiveLabel} 

\let\oldabs\abs
\def\abs{\@ifstar{\oldabs}{\oldabs*}}
\makenomenclature
\begin{document}


%----------------------------------------------------------------------------------------
%   Өөрийн мэдээллээ оруулах хэсэг
%----------------------------------------------------------------------------------------

% Дипломийн ажлын сэдэв
\title{ЗУРГАНД СУУРИЛСАН БАРАА ТАНИХ БОЛОН АУДИТЫН ШИЙДВЭР ГАРГАХ АВТОМАТ СИСТЕМ}
% Дипломын ажлын англи нэр
\titleEng{Image-Based Product Recognition and Automated Audit Decision System}
% Өөрийн овог нэрийг бүтнээр нь бичнэ
\author{Түвшинжаргалын Анар}
% Өөрийн овгийн эхний үсэг нэрээ бичнэ
\authorShort{Т.Анар}
% Удирдагчийн зэрэг цол овгийн эхний үсэг нэр
\supervisor{Др. Ч.Алтангэрэл}
% Хамтарсан удирдагчийн зэрэг цол овгийн эхний үсэг нэр
\cosupervisor{Др. Г.Амарсанаа}

% СиСи дугаар 
\sisiId{12D1SIT0001}
% Их сургуулийн нэр
\university{МОНГОЛ УЛСЫН ИХ СУРГУУЛЬ}
% Бүрэлдэхүүн сургуулийн нэр
\faculty{ХЭРЭГЛЭЭНИЙ ШИНЖЛЭХ УХААН, ИНЖЕНЕРЧЛЭЛИЙН СУРГУУЛЬ}
% Тэнхимийн нэр
\department{МЭДЭЭЛЭЛ, КОМПЬЮТЕРИЙН УХААНЫ ТЭНХИМ}
% Зэргийн нэр
\degreeName{Бакалаврын судалгааны ажил}
% Суралцаж буй хөтөлбөрийн нэр
\programeName{Мэдээллийн технологи (D061303)}
% Хэвлэгдсэн газар
\cityName{Улаанбаатар}
% Хэвлэгдсэн огноо
\gradyear{2024 оны 05 сар}


%----------------------------------------------------------------------------------------
%   Доорх хэсгийг өөрчлөх шаардлагагүй
%----------------------------------------------------------------------------------------
\include{main-pre}

% Удиртгалыг оруулж ирэх ба abstract.tex файлд удиртгалаа бичнэ
\begin{abstract}
Энэхүү судалгааны ажил нь жижиглэн худалдааны салбарт дэлгүүрийн лангуун дээрх бараа бүтээгдэхүүнийг зургаас автоматаар таньж, тооллого болон аудитын шийдвэр гаргалтыг хөнгөвчлөх системийг хөгжүүлэхэд оршино. Орчин үед дэлгүүрийн бараа материалын тооллого нь ихэвчлэн хүний оролцоотойгоор хийгдэж байгаа бөгөөд энэ нь цаг хугацаа их шаардсан, хөдөлмөрийн үр ашиггүй хэлбэрээс гадна алдаа гарах магадлал өндөр байдаг сул талтай. Бид энэхүү асуудлыг шийдвэрлэхийн тулд YOLO (You Only Look Once) архитектурын сүүлийн үеийн хувилбар болох YOLOv8 гүн сургалтын объектыг илрүүлэх загвар, FastAPI дээр суурилсан өндөр хурдны сүлжээний үйлчилгээ (Backend), болон хэрэглэгчид зориулсан олон платформт гар утасны Flutter аппликейшнээс бүрдсэн цогц мэдээллийн системийг санал болгож байна. Системийн ажиллах зарчмын хувьд гар утасны аппликейшнээр дамжуулан авсан лангууны зургийг сервер рүү илгээж, YOLOv8 загвар ашиглан зурган дээрх барааны байршил, төрлийг олж таних ба танигдсан үр дүнг өгөгдлийн санд урьдчилан тодорхойлсон хүлээгдэж буй барааны жагсаалттай харьцуулан автоматаар аудит хийх алгоритмыг хэрэгжүүлсэн. Туршилтын үр дүнгээс харахад сургагдсан загвар нь mAP50 хэмжүүрээр өндөр нарийвчлал үзүүлсэн ба нэг зургийг боловсруулах дундаж хугацаа нь үр дүнтэй (real-time) байлаа. Ингэснээр хөгжүүлсэн систем нь хүний оролцоог багасгаж, дэлгүүрийн менежментийн болон бараа нийлүүлэлтийн хяналт шалгалтын үр ашгийг эрс сайжруулах бүрэн боломжтойг батлан харуулж байна.
\end{abstract}


%----------------------------------------------------------------------------------------
%   Дипломын үндсэн хэсэг эндээс эхэлнэ
%----------------------------------------------------------------------------------------

\subfile{sections/introduction.tex}
\subfile{sections/related_work.tex}
\subfile{sections/methodology.tex}
\subfile{sections/implementation.tex}
\subfile{sections/evaluation.tex}

%----------------------------------------------------------------------------------------
%   Дүгнэлт эндээс эхэлнэ
%----------------------------------------------------------------------------------------
\subfile{sections/conclusion.tex}


%----------------------------------------------------------------------------------------
%   Дипломын номзүй, хавсралтын хэсэг эндээс эхэлнэ
%----------------------------------------------------------------------------------------

\singlespace
\addcontentsline{toc}{part}{НОМ ЗҮЙ}
\begin{thebibliography}{}
    \bibitem{yolov8_ultralytics}
    Jocher, G., Chaurasia, A., \& Qiu, J. (2023). YOLO by Ultralytics (Version 8.0.0) [Computer software]. \url{https://github.com/ultralytics/ultralytics}
    
    \bibitem{yolo_original}
    J. Redmon, S. Divvala, R. Girshick, and A. Farhadi, "You Only Look Once: Unified, Real-Time Object Detection", in \textit{Proceedings of the IEEE CVPR}, 2016, pp. 779-788.
    
    \bibitem{fastapi}
    S. Ramírez, "FastAPI Framework, high performance, easy to learn, fast to code", \url{https://fastapi.tiangolo.com/}, Accessed: 2024.
    
    \bibitem{flutter}
    Google, "Flutter - Build apps for any screen", \url{https://flutter.dev/}, Accessed: 2024.
    
    \bibitem{trax_retail}
    Trax Retail, "Computer Vision for Retail Execution", \url{https://traxretail.com/}, Accessed: 2024.
    
    \bibitem{amazon_go}
    Amazon, "Amazon Go and Just Walk Out technology", \url{https://www.amazon.com/b?node=16008589011}, Accessed: 2024.
\end{thebibliography}


%----------------------------------------------------------------------------------------
%   Хавсралтууд эндээс эхэлнэ
%----------------------------------------------------------------------------------------
\appendix
\addcontentsline{toc}{part}{ХАВСРАЛТ}

\subfile{appendix/code.tex}

\end{document}
